\section{Proposta de avaliação de impacto}
\label{sec:avaliacao}

A Tabela~\ref{tab:plano} resume o plano de avaliação.

\begin{table}[ht]
\centering\small
\caption{Plano de avaliação}
\label{tab:plano}
\begin{tabularx}{\textwidth}{@{}LLLL@{}}
\toprule
Pergunta & Sub-perguntas & Método & Indicadores e fonte \\
\midrule
A lei melhorou a saúde ao nascer? & Peso médio; baixo peso; duração da gestação; perda fetal & Diferenças em diferenças com estudo de evento, Chapada CE $\times$ RN & SINASC; SIM \\
A lei mudou a exposição de trabalhadores? & Intoxicações não intencionais sobem com a aplicação manual? & Idem; tentativa de suicídio como placebo & SINAN; SIH \\
A lei custou produção? & Área, produção e valor da banana; emprego formal & Idem & PAM; RAIS \\
Para quem? & O efeito é maior onde a agricultura é não familiar? & Diferença tripla & Censo Agropecuário \\
\bottomrule
\end{tabularx}
\end{table}

\subsection{Parâmetro}

Seja $Y_{it}(1)$ o desfecho do município $i$ no período $t$ sob a lei e
$Y_{it}(0)$ o desfecho sem ela. Só um dos dois é observado. O parâmetro de
interesse é o efeito médio sobre os tratados em cada período posterior à lei,
\[
ATT(t) = E\big[\,Y_{t}(1) - Y_{t}(0) \,\big|\, G = 2019\,\big], \qquad t \geq 2019,
\]
em que $G = 2019$ identifica os municípios cearenses da Chapada. É um efeito de
intenção de tratar no nível do município: o efeito de estar sob a lei, por
todos os caminhos da Seção~\ref{sec:caminhos}, e não o efeito de deixar de ser
pulverizado. A conversão para o efeito sobre a gestante exposta é exercício de
calibração, com a fração exposta declarada como hipótese.

\subsection{Estratégia de identificação}

O contrafactual é construído com os municípios potiguares da Chapada. Um grupo
de comparação válido precisa ter a mesma tendência, não ser afetado pelo
tratamento e reagir a ele da mesma forma \cite{gertler2016}. A identificação
apoia-se em quatro hipóteses.

\begin{enumerate}
\item \textbf{Tendências paralelas:} na ausência da lei, o desfecho dos
municípios cearenses da Chapada teria evoluído como o dos potiguares.
\item \textbf{Ausência de antecipação} antes do período de base.
\item \textbf{Sobreposição:} há municípios de controle comparáveis aos tratados.
Será examinada numa tabela de equilíbrio de linha de base, com população,
nascimentos, área de banana e participação da agricultura não familiar. O
controle inclui Mossoró, cidade-polo; as estimativas serão repetidas sem ela.
\item \textbf{O controle não é afetado pela lei}, e a composição das mães não
muda por causa dela.
\end{enumerate}

A Tabela~\ref{tab:ameacas} confronta cada ameaça com o que o desenho faz.

\begin{table}[ht]
\centering\small
\caption{Ameaças à validade interna}
\label{tab:ameacas}
\begin{tabularx}{\textwidth}{@{}p{3.1cm}LL@{}}
\toprule
Ameaça & O que é aqui & Resposta do desenho \\
\midrule
Perfil dos municípios & Municípios de fruticultura têm outro nível de desfecho & Fixo no tempo; sai na diferença. Covariáveis de linha de base \\
Seca e chuva & A seca de 2012--2017 afeta renda, saúde e a própria pulverização & Mesma chapada e mesmo clima nos dois lados \\
Choques comuns & Zika (2015--16), pandemia (2020--21) & Saem com o grupo de comparação \\
Choque estadual & Política cearense de saúde materna próxima de 2019 & \textbf{Não é fechada pela fronteira}: diferença tripla e desenho de calendário \\
Contaminação do controle & Deriva e aquífero cruzam a divisa; há relato de imprensa de aviões operando no lado potiguar até a divisa depois de 2019 & Viés contra achar efeito; sensibilidade sem os municípios colados à divisa \\
Antecipação & Projeto de lei de 2015, aprovação em 12/2018 & Período de base anterior a 2015 \\
Mudança no controle & Drones no lado potiguar a partir de cerca de 2021 & Sensibilidade com o pós-período até 2021 \\
Seleção para nascer vivo & Menos perda fetal altera quem nasce & Óbito fetal, razão de sexos e número de nascimentos como desfechos \\
Registro & Mudanças no formulário do SINASC; cobertura diferente por estado & Efeitos de período; conferência da cobertura nos dois estados \\
\bottomrule
\end{tabularx}
\end{table}

A fronteira deixa uma ameaça inteiramente aberta. Todos os tratados estão
no Ceará, de modo que qualquer mudança estadual simultânea se confunde com a
lei. A própria lei tem um segundo componente estadual: veda também a dispersão
aérea para controle de vetores em todo o Ceará. As especificações das Seções~\ref{sec:ddd} e \ref{sec:calendario} existem
para fechá-la.

\subsection{Estimador}

O estimador é o de \citeonline{callaway2021multiple}. Com uma única coorte tratada,
ele se reduz a comparações entre o grupo tratado e o grupo nunca tratado,
período a período, em relação a um período de base, e entrega um $ATT(t)$ para
cada período. Usa-se a versão duplamente robusta \cite{santanna2020drdid}, que
combina regressão do desfecho e escore de propensão sobre covariáveis de linha
de base e é consistente se um dos dois modelos estiver correto. As células são
ponderadas pelo número de nascimentos. Com dois ou três municípios tratados, o
ajuste por covariáveis tem pouco suporte; por isso reporta-se também a versão
sem covariáveis, que coincide com o estudo de evento por efeitos fixos de
município e de período, e a versão duplamente robusta é lida com essa
ressalva.

Três escolhas decorrem do grafo causal. As covariáveis são todas anteriores à
lei: população e urbanização em 2010 e participação da agricultura não familiar
em 2017.
Variáveis que a lei pode alterar, como cobertura de pré-natal e tipo de parto,
não entram, porque condicionar nelas fecharia parte do efeito. O período de base
é anterior ao projeto de lei de 2015, para acomodar antecipação. E a janela
termina em 19 de dezembro de 2024.

\paragraph{Previsão de forma.} A lei entra em vigor em janeiro de 2019. Uma
criança nascida em fevereiro teve a gestação quase toda sob o regime anterior.
O efeito deve, portanto, crescer ao longo de 2019 e estabilizar a partir das
concepções daquele ano, com perfil sazonal se a aplicação era sazonal. Um
efeito que surja inteiro em janeiro, sem rampa, conta contra a leitura causal.
A previsão fica registrada antes da estimação; dado o poder do desenho
(Seção~\ref{sec:poder}), serve como verificação descritiva, não como teste.

\subsection{Amostra}

\paragraph{Tratados.} Municípios cearenses da Chapada do Apodi. O núcleo são
Limoeiro do Norte e Quixeré, os dois com aplicação por aeronave medida no Censo
de 2006 (Tabela~\ref{tab:censo}) e sede do perímetro irrigado. A versão ampla
acrescenta Russas, citada pela literatura da região e sem registro de aeronave
em 2006, Jaguaruana, com um estabelecimento em 2006, e os demais municípios
cearenses da Chapada. A lista ampla será fechada antes da estimação.

\paragraph{Controles.} Os vinte e um municípios potiguares listados no projeto
de lei complementar que propôs a Região Integrada de Desenvolvimento da Chapada
do Apodi. A região não chegou a ser instituída; a lista serve como delimitação
externa ao projeto.

\paragraph{Período.} Na especificação principal, de 2010 a 19 de dezembro de
2024. A extensão do pré-período até 2005 depende da disponibilidade dos
arquivos nos dois estados e da forma de registro da idade gestacional nos anos
mais antigos, a conferir; ela importa para o cenário em que a variação de
exposição está em 2012 (Seção~\ref{sec:cenarios}).

\paragraph{Unidade.} Município de residência da mãe por trimestre de nascimento.
Não há amostragem: usa-se o universo dos nascimentos registrados, e a seleção
dos municípios é intencional. As listas de municípios são fixadas antes de se
examinar qualquer desfecho posterior a 2019 no recorte tratado contra
controle.

\subsection{Dados}

\begin{table}[ht]
\centering\small
\caption{Variáveis, fontes e situação}
\label{tab:dados}
\begin{tabularx}{\textwidth}{@{}p{3.3cm}LLp{2.2cm}@{}}
\toprule
Papel & Variável & Fonte & Situação \\
\midrule
Desfecho primário & Peso ao nascer (g) & SINASC & CE parcial; RN a baixar \\
Desfechos secundários & Baixo peso; semanas de gestação; prematuridade & SINASC & CE parcial; RN a baixar \\
Perda fetal & Óbito fetal; razão de sexos; número de nascimentos & SIM; SINASC & CE parcial; RN a baixar \\
Primeiro estágio & Área aplicada por avião, por município e mês & Relatórios mensais ao Ministério da Agricultura & Pedido de acesso a protocolar \\
Linha de base do método & Estabelecimentos com aplicação por aeronave, por tipo de agricultura & Censo Agropecuário 2006 & Obtido \\
Substituição & Modalidade de aplicação receitada & Receituário agronômico (ADAGRI) & Pedido a protocolar \\
Exposição ocupacional & Intoxicação não intencional por agrotóxico agrícola & SINAN; SIH & CE obtido; RN a baixar \\
Economia & Área, produção e valor da banana; vínculos formais & PAM; RAIS & PAM disponível; RAIS a obter \\
Composição & Idade, escolaridade e gemelaridade & SINASC & CE parcial; RN a baixar \\
Heterogeneidade & Participação da agricultura não familiar & Censo Agropecuário 2017 & Obtido \\
\bottomrule
\end{tabularx}
\end{table}

Os microdados de nascimento e de mortalidade são públicos e anonimizados. As
definições seguem o padrão da literatura: baixo peso abaixo de 2.500 gramas e
prematuridade abaixo de 37 semanas. A notificação de intoxicação mede também a
capacidade de notificar, que depende de vigilância e de rede hospitalar
\cite{rodrigues2022tese}. Por isso o projeto confere a evolução dessa
capacidade nos dois lados, usa a internação como segunda medida e trata a
tentativa de suicídio como desfecho placebo.

\subsection{Diferença tripla}
\label{sec:ddd}

Um choque estadual atinge o Ceará inteiro; a lei só importa onde havia avião.
Compara-se então a Chapada com o restante de cada estado:
\[
DDD = \big[(\text{Chapada} - \text{resto})_{CE}^{\text{pós}-\text{pré}}\big] -
      \big[(\text{Chapada} - \text{resto})_{RN}^{\text{pós}-\text{pré}}\big].
\]
A primeira diferença remove o que é cearense e varia no tempo. A segunda remove
o que é próprio de municípios de fruticultura irrigada e varia no tempo. O
restante de cada estado são os municípios sem aplicação por aeronave em 2006;
os poucos municípios com aeronave fora da Chapada ficam de fora dos dois
grupos. A identificação exige uma única hipótese de tendências paralelas, sobre
a diferença entre os dois contrastes \cite{olden2022triple}: um choque
estadual precisa atingir a Chapada e o restante do estado por igual.

\subsection{Desenho de calendário}
\label{sec:calendario}

O avião voava na quadra chuvosa, de fevereiro a maio, e
\citeonline{calzada2023bananas} encontram o efeito no primeiro e no segundo trimestre
de gestação. Dentro de cada município e ano há, portanto, gestações cujo início
atravessa a estação de aplicação e gestações que não. Compara-se a diferença
entre esses dois grupos, antes e depois da lei, no lado cearense e no potiguar.

A célula de comparação é o município por coorte de concepção, e não por mês de
nascimento, porque a prematuridade desloca o mês em que a criança nasce. A
comparação é feita dentro do município e absorve o que varia por município e
ano, o choque estadual inclusive. Fica aberta a possibilidade de a sazonalidade
dos nascimentos ter mudado em 2019 de modo diferente no lado cearense. A
composição sazonal das mães é ameaça conhecida \cite{buckles2013,
currie2013}, e o teste de composição se repete por estação. O ganho desse
desenho é de identificação; ele não aumenta o efeito esperado, e o seu poder
será calculado à parte.

\subsection{Heterogeneidade}

A lei atinge quem aplicava por avião, e o Censo de 2006 mostra que eram
sobretudo estabelecimentos não familiares. A previsão é que o efeito se
concentre onde a agricultura não familiar pesa mais. Testa-se com o corte da
amostra pela participação da agricultura não familiar e, na diferença tripla,
com essa participação no lugar da indicadora da Chapada. A medida é municipal e
não permite afirmações sobre famílias.

\subsection{O que o desenho estima em cada cenário do primeiro estágio}
\label{sec:cenarios}

Se o avião voava até 2018, a lei removeu uma exposição em curso e o contraste de
2019 a mede. Se a seca já o havia parado em 2012, a exposição não muda em 2019,
e o contraste informativo é o de 2012: o lado cearense deixa de pulverizar, e a
seca atinge os dois lados. A Tabela~\ref{tab:cenarios} registra a regra de
decisão antes de o registro chegar.

\begin{table}[ht]
\centering\small
\caption{Regra de decisão conforme o registro da aplicação aérea}
\label{tab:cenarios}
\begin{tabularx}{\textwidth}{@{}LLL@{}}
\toprule
O que o registro mostrar & Contraste principal & O que se estima \\
\midrule
Voo no lado cearense até 2018 & 2019 & Efeito de a lei remover a exposição \\
Lado cearense para em 2012; lado potiguar não voava ou segue voando & 2012, com 2019 como teste de que a lei nada acrescentou & Efeito de cessar a exposição; a lei consolidou \\
Os dois lados param em 2012 & Calendário anterior a 2012, no lado cearense & Efeito da exposição sazonal \\
\bottomrule
\end{tabularx}
\end{table}

No segundo cenário a identificação pede uma hipótese adicional e forte: a seca
afetou renda e saúde igualmente nos dois lados e só de um lado desligou o
avião. Os períodos anteriores a 2012 passam a ser os de teste de tendência, e
os de 2012 a 2018 deixam de sê-lo. No terceiro cenário não há grupo de
comparação, e a diferença entre estações só se lê como efeito da exposição sob
a hipótese de que a sazonalidade do peso ao nascer, sem o avião, seria a
observada em municípios vizinhos. Estimador, data de corte e poder de cada
cenário serão especificados antes de o registro ser aberto.

\subsection{Inferência}

Com dois a nove municípios tratados, o erro-padrão agrupado convencional não é
confiável. A inferência usa os métodos de \citeonline{conley2011} e
\citeonline{ferman2019few}, que obtêm a distribuição do estimador a partir dos
controles, e testes de permutação entre municípios de controle. Com vinte e um
controles, o menor valor-p que a permutação alcança é da ordem de 5\%, e esses
métodos têm baixo poder. A sensibilidade a violações de tendências paralelas
segue \citeonline{rambachan2023credible}. Os desfechos secundários recebem a
correção de Holm para comparações múltiplas.

\subsection{Validade externa}

O resultado vale para a proibição de um método de aplicação em fruticultura
irrigada de exportação, com fungicidas em banana, numa região semiárida. Não se
transporta para herbicidas em grãos. Interessa, porém, a uma decisão que outros
estados podem tomar: o Supremo reconheceu a competência estadual para legislar
sobre a matéria.
